\documentclass[a4paper]{article}

% Español (México) con XeLaTeX: fontspec para fuentes Unicode, polyglossia
% para la división silábica y los nombres del documento en la variante mexicana.
\usepackage{fontspec}
\usepackage{polyglossia}
\setdefaultlanguage[variant=mexican]{spanish}
% Los nombres chinos van como «pinyin (original)»: xeCJK da a los caracteres
% Han su propia fuente; sin ella XeLaTeX los omite con un aviso.
% FandolSong no cubre algunos caracteres de nombres (U+73FA); WenQuanYi Zen Hei sí.
\usepackage[AutoFallBack=true]{xeCJK}
\setCJKmainfont{FandolSong-Regular.otf}
\setCJKfallbackfamilyfont{\CJKrmdefault}{WenQuanYi Zen Hei}
% polyglossia no traduce los nombres de listings.
\AtBeginDocument{\providecommand{\lstlistingname}{}\renewcommand{\lstlistingname}{Listado}%
  \providecommand{\lstlistlistingname}{}\renewcommand{\lstlistlistingname}{Índice de listados}}
\usepackage{amsmath,amssymb}
\usepackage{graphicx}
\usepackage[margin=2.5cm]{geometry}
\usepackage[most]{tcolorbox}
\usepackage{etoolbox}
\usepackage{subcaption}
\usepackage{listings}
\usepackage{hyperref}
\usepackage{booktabs}
\usepackage{float}

\newtcolorbox{knowledgebox}[1]{
    enhanced, colback=blue!5!white, colframe=blue!75!black, colbacktitle=blue!75!black,
    coltitle=white, fonttitle=\bfseries, title={#1}, boxrule=1pt, sharp corners
}

\newtcolorbox{importantbox}[1]{
    enhanced, colback=yellow!10!white, colframe=yellow!80!black, colbacktitle=yellow!80!black,
    coltitle=black, fonttitle=\bfseries, title={#1}, sharp corners
}

\newtcolorbox{warningbox}[1]{
    enhanced, colback=red!5!white, colframe=red!75!black, colbacktitle=red!75!black,
    coltitle=white, fonttitle=\bfseries, title={#1}, sharp corners
}

\lstset{
    basicstyle=\ttfamily\small,
    keywordstyle=\color{blue},
    stringstyle=\color{red!60!black},
    commentstyle=\color{green!60!black},
    breaklines=true,
    frame=single,
    numbers=left,
    numberstyle=\tiny\color{gray},
    captionpos=b,
    extendedchars=true
}

\newcommand{\notetitle}{CS224R Lecture 11: Aprendizaje por refuerzo basado en modelos}
\newcommand{\noteauthors}{Elaboradas a partir de materiales públicos del curso}
\newcommand{\notedate}{Primavera de 2025}
\newcommand{\videochannel}{Stanford Online}
\newcommand{\videopublishdate}{2025}
\newcommand{\videoduration}{Aproximadamente 80 minutos}
\newcommand{\videourl}{https://cs224r.stanford.edu/}
\newcommand{\videocoverpath}{cover.jpg}
\newcommand{\repourl}{https://github.com/hqhq1025/ai-course-notes}


% Footer with repo info
\usepackage{fancyhdr}
\pagestyle{fancy}
\fancyfoot[C]{\thepage}
\fancyhf{}
\fancyfoot[R]{\footnotesize\color{gray}\href{\repourl}{AI Course Notes}}
\renewcommand{\headrulewidth}{0pt}
\renewcommand{\footrulewidth}{0pt}

\begin{document}
\begin{titlepage}
\centering
{\Large Notas del curso\par}
\vspace{1.2cm}
{\huge\bfseries \notetitle\par}
\vspace{0.8cm}
{\large \noteauthors\par}
\vspace{0.3cm}
{\large \notedate\par}
\vspace{1.1cm}

\ifdefempty{\videocoverpath}{
{\small No se proporcionó imagen de portada.\par}
}{
\includegraphics[width=0.82\textwidth,height=0.45\textheight,keepaspectratio]{\videocoverpath}\par
}

\vfill
\begin{tcolorbox}[width=0.9\textwidth, colback=black!2!white, colframe=black!60, sharp corners]
\textbf{Autor o canal del video}: \videochannel\par
\textbf{Fecha de publicación}: \videopublishdate\par
\textbf{Duración del video}: \videoduration\par
\textbf{Ponente}: Chelsea Finn\par
\textbf{Página del curso}: \href{https://cs224r.stanford.edu/}{cs224r.stanford.edu}\par
\textbf{Página del proyecto}: \href{\repourl}{\nolinkurl{\repourl}}\par
\end{tcolorbox}
\end{titlepage}

\tableofcontents
\newpage

%% --- Inicio del contenido --- %%
\section{Repaso del curso: panorama de la clasificación de algoritmos}

Al inicio de esta clase, Chelsea Finn hizo un repaso de alto nivel de todos los algoritmos que se han visto en el curso hasta ahora.

\begin{importantbox}{Marco de clasificación de algoritmos}
\begin{itemize}
    \item \textbf{En línea (Online) vs. fuera de línea (Offline)}: los algoritmos en línea permiten interactuar de forma continua con el entorno durante el entrenamiento para recopilar datos; los algoritmos fuera de línea solo usan un conjunto de datos fijo y dado.
    \item \textbf{Dentro de la categoría Online}:
    \begin{itemize}
        \item On-policy RL: solo usa datos generados por la política actual (por ejemplo, REINFORCE, PPO)
        \item Off-policy RL: puede usar datos de políticas anteriores (por ejemplo, SAC, TD3)
    \end{itemize}
    \item \textbf{Dentro de la categoría Offline}:
    \begin{itemize}
        \item Imitation Learning: clonación de comportamiento (BC), DAgger
        \item Offline RL: CQL y otros métodos que restringen la política para que no se aleje de la distribución de los datos
    \end{itemize}
\end{itemize}
\end{importantbox}

\section{Idea central de Model-Based RL}

\subsection{Qué es Model-Based RL}

A diferencia de los métodos model-free, model-based RL aprende de forma explícita un \textbf{modelo de dinámica} (dynamics model) del entorno y después lo usa para tomar decisiones.

\begin{importantbox}{Modelo de dinámica}
Se aprende una función $f_\phi(\mathbf{s}, \mathbf{a})$ que aproxima la transición de estados del entorno $p(\mathbf{s}' | \mathbf{s}, \mathbf{a})$:
$$f_\phi(\mathbf{s}, \mathbf{a}) \approx \mathbf{s}'$$
El objetivo del entrenamiento es minimizar el error de predicción:
$$\min_\phi \sum_i \|f_\phi(\mathbf{s}_i, \mathbf{a}_i) - \mathbf{s}_i'\|^2$$
donde los datos $\mathcal{D} = \{(\mathbf{s}, \mathbf{a}, \mathbf{s}')_i\}$ provienen de la interacción con el entorno.
\end{importantbox}

\begin{figure}[H]
\centering
\includegraphics[width=0.9\textwidth]{slides-images/slide-05.jpg}
\caption{Temario de esta clase: contenidos centrales de Model-Based RL\protect\footnotemark}
\end{figure}
\footnotetext{Slides, página 5.}

\subsection{Por qué aprender un modelo}

\begin{knowledgebox}{Disyuntiva central entre Model-Based y Model-Free}
\begin{itemize}
    \item \textbf{Eficiencia de datos}: los métodos model-based suelen ser mucho más eficientes en el uso de datos que los métodos model-free. El modelo aprendido puede usarse para «imaginar» un sinnúmero de trayectorias posibles sin necesidad de ejecutarlas en la realidad.
    \item \textbf{Costo}: el modelo no es perfecto, y las decisiones basadas en un modelo inexacto pueden degradar el desempeño. Este es el problema del \textbf{sesgo del modelo} (model bias).
\end{itemize}
\end{knowledgebox}

\subsection{Resumen de la sección}

La idea central de Model-Based RL es «primero aprender un modelo del mundo y después usar ese modelo para planificar». Resulta especialmente valioso cuando los datos son escasos (por ejemplo, en robótica), pero exige manejar con cuidado el error del modelo.

\section{Cómo usar el modelo aprendido}

\subsection{Método 1: planificación basada en el modelo (Planning)}

Con un modelo de dinámica disponible, se puede buscar directamente en el modelo la secuencia de acciones óptima:

$$\max_{\mathbf{a}_{t:t+H}} \sum_t r(\mathbf{s}_t, \mathbf{a}_t)$$

\begin{figure}[H]
\centering
\includegraphics[width=0.9\textwidth]{slides-images/slide-15.jpg}
\caption{Método 1b: planificación mediante optimización sin gradientes por muestreo\protect\footnotemark}
\end{figure}
\footnotetext{Slides, página 15.}

Flujo concreto del algoritmo:
\begin{enumerate}
    \item Ejecutar alguna política (por ejemplo, una política aleatoria) para recolectar datos $\mathcal{D} = \{(\mathbf{s}, \mathbf{a}, \mathbf{s}')_i\}$
    \item Aprender el modelo $f_\phi(\mathbf{s}, \mathbf{a})$
    \item Muestrear iterativamente secuencias de acciones y elegir la mejor acción mediante la simulación hacia adelante con el modelo
\end{enumerate}

\subsubsection{Método de disparo aleatorio (Random Shooting)}

El método más sencillo consiste en muestrear aleatoriamente $N$ secuencias de acciones, hacer el rollout de cada secuencia en el modelo y elegir la primera acción de la secuencia con la mayor recompensa acumulada.

\subsubsection{CEM (Cross-Entropy Method)}

CEM es un método de optimización sin gradientes más avanzado:
\begin{enumerate}
    \item Inicializar una distribución de acciones (por ejemplo, una distribución gaussiana)
    \item Muestrear $N$ secuencias de acciones de la distribución
    \item Evaluar en el modelo la recompensa acumulada de cada secuencia
    \item Seleccionar las top-$K$ secuencias «de élite»
    \item Actualizar la media y la varianza de la distribución de acciones con las secuencias de élite
    \item Repetir durante varias iteraciones
\end{enumerate}

\begin{knowledgebox}{MPC (Model Predictive Control)}
En la ejecución real se suele adoptar el paradigma MPC: en cada paso solo se ejecuta la primera acción obtenida por la planificación y luego se vuelve a planificar desde el nuevo estado. Esto mitiga en parte el error acumulado del modelo en horizontes temporales largos.
\end{knowledgebox}

\subsubsection{Método 1a: planificación basada en gradientes}

Si el modelo es diferenciable, se puede optimizar calculando directamente el gradiente respecto a la secuencia de acciones. Sin embargo, en la práctica este método tiende a quedar atrapado en óptimos locales y exige que el modelo sea suficientemente suave.

\subsection{Método 2: usar el modelo para generar datos sintéticos}

\begin{figure}[H]
\centering
\includegraphics[width=0.9\textwidth]{slides-images/slide-25.jpg}
\caption{Dos vías de optimización de políticas Model-Based\protect\footnotemark}
\end{figure}
\footnotetext{Slides, página 25.}

Principales limitaciones de los métodos de planificación:
\begin{itemize}
    \item Alto costo de cómputo en tiempo de prueba (hay que optimizar en cada paso)
    \item Están limitados por el problema del horizonte corto
\end{itemize}

Dos enfoques para resolver el problema del horizonte largo:
\begin{enumerate}
    \item \textbf{Planificación con función de valor terminal}: planificación de horizonte corto + una función de valor aprendida como estimación del valor terminal
    \item \textbf{Métodos estilo Dyna}: usar el modelo para generar datos sintéticos que refuercen el entrenamiento de RL model-free
\end{enumerate}

\begin{importantbox}{Algoritmo Dyna}
Dyna es un paradigma clásico de RL model-based:
\begin{enumerate}
    \item Recolectar datos reales y entrenar o actualizar el modelo de dinámica
    \item Usar el modelo para generar datos sintéticos (synthetic rollouts)
    \item Mezclar los datos sintéticos con los reales y actualizar la política con un algoritmo de RL model-free
\end{enumerate}
Su ventaja central es aprovechar la capacidad de generalización del modelo para «amplificar» una cantidad limitada de datos reales en una gran cantidad de datos de entrenamiento.
\end{importantbox}

\subsection{Manejo del error del modelo}

\begin{warningbox}{Efecto acumulativo del error del modelo}
Los pequeños errores del modelo en cada paso se acumulan a lo largo del eje temporal. Si el horizonte de planificación es $H$ y el error por paso es $\epsilon$, el error total después de $H$ pasos puede llegar a $O(H \cdot \epsilon)$ o ser aún mayor (si la dinámica es inestable). Por ello, confiar ciegamente en el modelo lleva a que la política «explote» los defectos del modelo (model exploitation).
\end{warningbox}

Estrategias de mitigación:
\begin{itemize}
    \item \textbf{Ensamble de modelos (Ensemble)}: entrenar varios modelos y usar las predicciones de los distintos modelos para estimar la incertidumbre
    \item \textbf{Rollout de horizonte corto}: limitar el número de pasos del rollout en el modelo
    \item \textbf{Actualización continua del modelo}: actualizar el modelo constantemente a medida que se recolectan datos nuevos
\end{itemize}

\subsection{Resumen de la sección}

El modelo puede usarse de dos maneras: (1) planificación directa (optimizar la secuencia de acciones) y (2) generación de datos sintéticos para reforzar el RL model-free. Ambas requieren manejar adecuadamente el error del modelo.
\section{Estudio de caso: Model-Based RL en manipulación diestra}

El ponente presentó un caso concreto en el que se usa model-based RL para que una mano robótica diestra manipule objetos.

\begin{figure}[H]
\centering
\includegraphics[width=0.9\textwidth]{slides-images/slide-35.jpg}
\caption{Resultados de los experimentos de ablación en simulación para la manipulación diestra\protect\footnotemark}
\end{figure}
\footnotetext{Slides, página 35.}

Hallazgos experimentales principales:
\begin{itemize}
    \item Se necesita una arquitectura de modelo suficientemente grande (2x250 o 2x500 unidades ocultas)
    \item Se requieren al menos 3 miembros del ensamble para obtener estimaciones de incertidumbre confiables
    \item El horizonte de planificación implica un compromiso: si es demasiado corto, no se ve lo bastante lejos en el futuro; si es demasiado largo, se acumula el error del modelo
    \item El controlador CEM mejorado (por ejemplo, Filtering + Reward Weighting) es mucho mejor que Random Shooting
\end{itemize}

\subsection{Cuándo usar Model-Based RL}

\begin{importantbox}{Escenarios de aplicación de Model-Based RL}
\begin{itemize}
    \item La recolección de datos es costosa (por ejemplo, con robots reales)
    \item El espacio de estados es de dimensión relativamente baja y predecible
    \item Se necesita adaptarse con rapidez a tareas nuevas (el modelo es transferible)
    \item La dinámica tiene cierta estructura que se puede aprender
\end{itemize}
Escenarios no adecuados: entornos con observaciones de alta dimensión (como imágenes) y dinámica altamente estocástica.
\end{importantbox}

\subsection{Resumen de la sección}

El caso de manipulación diestra muestra la ventaja de model-based RL en eficiencia de datos, así como la influencia decisiva de hiperparámetros como el ensamble y el horizonte de planificación en el desempeño.
\section{Síntesis y ampliación}

El Model-Based RL es un componente importante de la caja de herramientas del RL; su principal ventaja es la eficiencia en el uso de datos.

Puntos principales:
\begin{enumerate}
    \item Aprender un modelo de la dinámica puede reducir mucho los datos de interacción real que se necesitan
    \item El modelo puede usarse para planning (CEM/MPC) o para generar datos sintéticos (al estilo de Dyna)
    \item El error del modelo es el desafío central; se mitiga con ensambles, rollouts de horizonte corto y actualización continua
    \item Es especialmente valioso en escenarios donde los datos son costosos, como la manipulación robótica
\end{enumerate}

\begin{knowledgebox}{Lecturas adicionales}
\begin{itemize}
    \item Nagabandi et al., ``Neural Network Dynamics for Model-Based Deep RL with Model-Free Fine-Tuning,'' ICRA 2018
    \item Chua et al., ``Deep Reinforcement Learning in a Handful of Trials using Probabilistic Dynamics Models (PETS),'' NeurIPS 2018
    \item Janner et al., ``When to Trust Your Model: Model-Based Policy Optimization (MBPO),'' NeurIPS 2019
    \item Sutton, ``Dyna, an Integrated Architecture for Learning, Planning, and Reacting,'' 1991
\end{itemize}
\end{knowledgebox}

\end{document}
